%==============================================================================
% PARTIDA 03: CANALIZACIONES, TUBERIAS Y CONDUCTOS
% Codigos: OE.5.2.2.1 al OE.5.2.2.7
%==============================================================================
\subsection{OE.5.2.2.1 Tuberia PVC Liviana (PV-L) 3/4'' - Canalizacion de Alumbrado}

\subsubsection{Especificaciones Tecnicas}

La tuberia PVC liviana (PV-L) de 3/4'' de diametro se utilizara para la canalizacion de los circuitos de alumbrado (C1, C3, C5). Debe ser de policloruro de vinilo (PVC) tipo liviano, segun la NTP 399.006, de color naranja o crema, con espesor de pared de 1.5 mm minimo. La tuberia se instalara empotrada en techos, losas y muros.

La canalizacion debe ser continua, sin interrupciones, desde la caja de derivacion hasta la caja de salida de alumbrado. No se permiten empalmes de conductores dentro de la tuberia.

\subsubsection{Requisitos de Materiales}

\begin{itemize}
  \item Tuberia PVC liviana (PV-L) de 3/4'' de diametro, longitud de 3 m por unidad.
  \item Curvas PVC livianas de 3/4'' (45° y 90°) para cambios de direccion.
  \item Uniones universales de PVC de 3/4'' para empalmes de tuberia.
  \item Pegamento PVC (cemento solvente) para fijacion de uniones y accesorios.
  \item Cinta teflon para sellado de roscas.
  \item Abrazaderas tipo gancho para fijacion temporal durante el vaciado.
\end{itemize}

\subsubsection{Metodo de Ejecucion}

\begin{enumerate}
  \item Trazar la ruta de la canalizacion en techos y muros segun el plano de alumbrado (IE-02, IE-03, IE-04).
  \item Cortar la tuberia PVC a las longitudes requeridas utilizando segueta o cortatubos.
  \item Aplicar pegamento PVC en los extremos de la tuberia y en el interior de las uniones o curvas.
  \item Ensamblar la tuberia con las curvas y uniones, asegurando una penetracion completa (2/3 de la longitud del accesorio).
  \item Fijar la tuberia a la estructura mediante abrazaderas temporales, con una separacion maxima de 1.00 m entre puntos de fijacion.
  \item Verificar que la tuberia quede continua y sin obstrucciones realizando una prueba de paso con una guia de alambre.
  \item Empotrar la tuberia en la losa o muro durante el vaciado de concreto o tarrajeo.
\end{enumerate}

\subsubsection{Control de Calidad}

\begin{itemize}
  \item Verificar que no existan aplastamientos o deformaciones en la tuberia.
  \item Comprobar que las curvas no presenten reduccion de seccion mayor al 10\%.
  \item Realizar prueba de paso con guia de alambre en cada tramo de tuberia.
  \item Inspeccionar que todas las uniones esten selladas con pegamento PVC.
  \item Verificar que la separacion entre abrazaderas no exceda 1.00 m.
\end{itemize}

\subsubsection{Metodo de Medicion}

Se medira por metro lineal (m) de tuberia PVC liviana de 3/4'' instalada, incluyendo curvas, uniones y accesorios necesarios para la canalizacion completa.

\subsubsection{Unidad de Medida}

m (metro lineal).

\subsubsection{Forma de Pago}

Pago por metro lineal instalado y verificado mediante prueba de paso, aprobado por supervision.

%------------------------------------------------------------------------------
\subsection{OE.5.2.2.2 Tuberia PVC Pesada (PV-P) 3/4'' - Canalizacion de Tomacorrientes}

\subsubsection{Especificaciones Tecnicas}

La tuberia PVC pesada (PV-P) de 3/4'' de diametro se utiliza para la canalizacion de los circuitos de tomacorrientes (C2, C4, C6). Debe ser de PVC tipo pesado, segun la NTP 399.006, de color naranja, con espesor de pared de 2.0 mm minimo, para brindar mayor proteccion mecanica a los conductores de los circuitos de tomacorrientes que tienen mayor carga y uso frecuente.

\subsubsection{Requisitos de Materiales}

\begin{itemize}
  \item Tuberia PVC pesada (PV-P) de 3/4'' de diametro, longitud de 3 m por unidad.
  \item Curvas PVC pesadas de 3/4'' (45° y 90°).
  \item Uniones universales de PVC pesado de 3/4''.
  \item Pegamento PVC (cemento solvente).
  \item Adaptadores de PVC de 3/4'' a rosca para conexion a cajas metalicas.
\end{itemize}

\subsubsection{Metodo de Ejecucion}

\begin{enumerate}
  \item Trazar la ruta de canalizacion en muros y losas segun planos IE-02, IE-03, IE-04.
  \item Cortar, ensamblar y pegar la tuberia con curvas y uniones siguiendo el mismo procedimiento de la tuberia liviana.
  \item Conectar la tuberia a las cajas rectangulares mediante adaptadores de PVC a rosca o conectores tipo "romex".
  \item Fijar la tuberia con una separacion maxima de 0.80 m entre puntos de soporte.
  \item Realizar prueba de paso con guia de alambre en cada tramo.
\end{enumerate}

\subsubsection{Control de Calidad}

\begin{itemize}
  \item Verificar el espesor de pared (minimo 2.0 mm) midiendo con calibrador.
  \item Comprobar que la tuberia no presente fisuras o roturas.
  \item Realizar prueba de paso en cada tramo.
  \item Inspeccionar la correcta conexion a las cajas mediante adaptadores.
\end{itemize}

\subsubsection{Metodo de Medicion}

Se medira por metro lineal (m) de tuberia PVC pesada de 3/4'' instalada.

\subsubsection{Unidad de Medida}

m (metro lineal).

\subsubsection{Forma de Pago}

Pago por metro lineal instalado y aprobado por supervision.

%------------------------------------------------------------------------------
\subsection{OE.5.2.2.3 Tuberia PVC Pesada (PV-P) 1'' - Canalizacion Alimentador General}

\subsubsection{Especificaciones Tecnicas}

La tuberia PVC pesada (PV-P) de 1'' de diametro se utiliza para la canalizacion del alimentador general, que transporta los conductores de 10 mm2 desde la caja portamedidor hasta el tablero general TG-01. El mayor diametro permite el alojamiento de conductores de mayor seccion y facilita el cableado.

\subsubsection{Requisitos de Materiales}

\begin{itemize}
  \item Tuberia PVC pesada (PV-P) de 1'' de diametro.
  \item Curvas PVC pesadas de 1'' (45° y 90°).
  \item Uniones universales de PVC pesado de 1''.
  \item Pegamento PVC.
  \item Adaptadores de 1'' a rosca.
\end{itemize}

\subsubsection{Metodo de Ejecucion}

\begin{enumerate}
  \item Trazar la ruta desde la caja portamedidor hasta TG-01 segun plano IE-01.
  \item Instalar la tuberia de 1'' siguiendo el mismo procedimiento de canalizacion general, con separacion maxima de soportes de 0.80 m.
  \item Asegurar que la tuberia tenga pendiente minima del 0.5\% hacia la caja portamedidor para evitar acumulacion de humedad.
  \item Realizar prueba de paso con guia de alambre.
\end{enumerate}

\subsubsection{Control de Calidad}

\begin{itemize}
  \item Verificar el diametro interior libre de obstrucciones.
  \item Comprobar la pendiente de la tuberia.
  \item Realizar prueba de paso con conductor de 10 mm2 de prueba.
\end{itemize}

\subsubsection{Metodo de Medicion}

Se medira por metro lineal (m) de tuberia PVC pesada de 1'' instalada.

\subsubsection{Unidad de Medida}

m (metro lineal).

\subsubsection{Forma de Pago}

Pago por metro lineal instalado y aprobado por supervision.

%------------------------------------------------------------------------------
\subsection{OE.5.2.2.4 a OE.5.2.2.7 - Curvas y Accesorios de Canalizacion}

\subsubsection{Especificaciones Tecnicas}

Las curvas PVC permiten cambios de direccion en las canalizaciones, manteniendo la seccion interior libre para el paso de conductores. Los accesorios de canalizacion incluyen uniones universales, adaptadores, conectores, pegamento PVC y otros elementos necesarios para completar la instalacion de tuberias.

Las curvas deben ser del mismo tipo y diametro que la tuberia correspondiente: curvas PV-L de 3/4'' para alumbrado, curvas PV-P de 3/4'' para tomacorrientes y curvas PV-P de 1'' para alimentador general.

\subsubsection{Requisitos de Materiales}

\begin{itemize}
  \item Curvas PVC liviana (PV-L) de 3/4'' en angulos de 45° y 90°.
  \item Curvas PVC pesada (PV-P) de 3/4'' en angulos de 45° y 90°.
  \item Curvas PVC pesada (PV-P) de 1'' en angulos de 45° y 90°.
  \item Uniones universales PVC de 3/4'' y 1''.
  \item Adaptadores de PVC a rosca de 3/4'' y 1''.
  \item Pegamento PVC (cemento solvente) en presentacion de 1/4 galon.
  \item Cinta teflon para sellado de conexiones roscadas.
\end{itemize}

\subsubsection{Metodo de Ejecucion}

\begin{enumerate}
  \item Seleccionar la curva o accesorio del tipo y diametro adecuado a la tuberia.
  \item Aplicar pegamento PVC en el extremo de la tuberia y en el interior del accesorio.
  \item Insertar la tuberia en el accesorio hasta el tope, girando ligeramente para distribuir el pegamento.
  \item Mantener la presion durante 15 segundos para asegurar la fijacion.
  \item Limpiar el exceso de pegamento con un trapo humedo.
  \item Para los accesorios de canalizacion (pegamento, uniones, conectores), se considera un global (glb) que incluye todos los insumos necesarios para completar la instalacion de tuberias segun los metrados.
\end{enumerate}

\subsubsection{Control de Calidad}

\begin{itemize}
  \item Verificar que las curvas no presenten deformaciones o reducciones de seccion.
  \item Comprobar que el angulo de la curva permita el paso libre de conductores.
  \item Inspeccionar que las uniones esten completamente selladas.
  \item Verificar que no existan rebabas en los bordes de los accesorios.
\end{itemize}

\subsubsection{Metodo de Medicion}

Las curvas se mediran por pieza (pza) instalada. Los accesorios de canalizacion (pegamento, uniones, conectores) se mediran como un global (glb) por toda la obra.

\subsubsection{Unidad de Medida}

pza (pieza) para curvas; glb (global) para accesorios.

\subsubsection{Forma de Pago}

Pago por pieza de curva instalada o por global de accesorios completado y verificado, aprobado por supervision.
